\subsubsection{Two infinite families with real rank-one singular sets}
\label{sec:orthq-families}

An unused block of weight two gives two explicit nonharmonic
families in $\R^{8k}$. Their singular set is the real rank-one
matrix cone, and their area-minimizing range has the sharp threshold
$k\ge3$.

Identify $\R^{4k}$ with $\mathbb H^k$ and let $J_1,J_2,J_3$ be
componentwise right multiplication by the quaternionic units
$\mathbf i,\mathbf j,\mathbf k$. For $x,y\in\mathbb H^k$, $k\ge2$, put
$R=|x|^2+|y|^2$ and $W=|x|^2|y|^2-\langle x,y\rangle^2$, and define
\begin{align}
 F_{k,2}&=(4k-3)
  (\langle x\mathbf i,y\rangle^2+\langle x\mathbf j,y\rangle^2)-W,
  \label{eq:orthq-two-families}\\
 F_{k,3}&=(4k-3)
  (\langle x\mathbf i,y\rangle^2+\langle x\mathbf j,y\rangle^2
                         +\langle x\mathbf k,y\rangle^2)-2W.
  \nonumber
\end{align}
Here the displayed pairings are real Euclidean inner products.
For $h=\sum_\alpha\bar x_\alpha y_\alpha$, one has
$\langle x\mathbf i,y\rangle=\Im_{\mathbf i}h$ and similarly for
the other units. Thus the second family is also the Gram composition
\begin{equation}\label{eq:orthq-quaternionic-Gram}
 F_{k,3}=(4k-3)|h|^2-(4k-5)(\Re h)^2-2|x|^2|y|^2.
\end{equation}

\begin{theorem}[Two families and their sharp threshold]
\label{thm:orthq-infinite-families}
For $k\ge2$ and $m=2,3$, the polynomial $F_{k,m}$ defines an
irreducible minimal hypercone $C_{k,m}\subset\R^{8k}$ with the
algebra and singular geometry in \cref{prop:orthq-construction}.
The cone is area minimizing if and only if $k\ge3$.
For $k\ge3$, the degree-eight potential
\begin{equation}\label{eq:orthq-potential}
 G_{k,m}=WF_{k,m}
\end{equation}
subcalibrates the full boundary current, which is strictly area
minimizing. The regular cone satisfies
\begin{equation}\label{eq:orthq-Hardy}
 Q_{C_{k,m}}(\varphi)\ge
 \left(16k^2-36k-\frac{31}{4}\right)
 \int_{C_{k,m}}r^{-2}\varphi^2.
\end{equation}
At $k=2$ both cones are unstable on their regular parts.
\end{theorem}

Both sequences arise from the following construction, which gives
their algebra, irreducibility, and transverse singular geometry.

Let $J_1,\ldots,J_m$ be skew-symmetric orthogonal endomorphisms of
$\R^\ell$ satisfying
\[
 J_iJ_j+J_jJ_i=-2\delta_{ij}\Id,
 \qquad m\ge2,\qquad\ell-m-1\ge2.
\]
For $x,y\in\R^\ell$, put
\begin{gather}
 R=|x|^2+|y|^2,\qquad
 W=|x|^2|y|^2-\langle x,y\rangle^2,\qquad
 Q=\sum_{i=1}^m\langle J_ix,y\rangle^2,
 \label{eq:orthq-data}\\
 \begin{gathered}
 a=\ell-m-2,\qquad b=m-1,\qquad d=a-b=\ell-2m-1,\\
 F=(\ell-3)Q-bW.
 \end{gathered}
 \label{eq:orthq-polynomial}
\end{gather}
The vectors $x,J_1x,\ldots,J_mx$ are mutually orthogonal and have
length $|x|$, so $0\le Q\le W\le R^2/4$.

\begin{proposition}[Construction and exact algebra]
\label{prop:orthq-construction}
The polynomial $F$ is real irreducible and nonharmonic, has nonempty
regular zero set, and has essential space $\R^{2\ell}$.
Its regular zero set is minimal, and
\begin{equation}\label{eq:orthq-minimality}
 |\nabla F|^2=4R(dF+abW),\qquad
 \cL(F)=-8F\bigl(2dF+2abW+abR^2\bigr).
\end{equation}
Moreover,
\[
 \cA_R(F)=\R[R,W,F]
\]
is a polynomial algebra on three generators of degrees $(2,4,4)$.
Its critical set is the geometric singular set of $C=\{F=0\}$:
\begin{equation}\label{eq:orthq-singular}
 \Sing C=\{W=0\}
 =\{(x,y):\operatorname{rank}_{\R}[x\ y]\le1\}.
\end{equation}
This set has dimension $\ell+1$ and is smooth away from the origin.
At each nonzero singular point the transverse tangent cone is the
Lawson quadratic cone of signature $(m,\ell-m-1)$.
\end{proposition}

\begin{proof}
The symmetric Clifford system associated with
$\Id,J_1,\ldots,J_m$ has quadratic coordinates
$|x|^2-|y|^2$, $2\langle x,y\rangle$, and
$2\langle J_ix,y\rangle$; see \cite[Section~3]{FKM81}.
The blocks
\begin{equation}\label{eq:orthq-blocks}
 E_0=R^2-4W,\qquad E_1=4Q,\qquad E_2=4(W-Q)
\end{equation}
satisfy \eqref{eq:orthogonal-quartic-data} with weights
$(2,m,\ell-m-1)$ by \cref{lem:Clifford-block-quartics}. In the normalization of
\eqref{eq:quartic-general-generators}, $T=W$ and
$F_{\rm gen}=F/(\ell-3)$. Thus \eqref{bh:reduced-table} gives
\begin{equation}\label{eq:orthq-table}
 \begin{gathered}
 (\Delta R,\Delta W,\Delta F)
 =(4\ell,2(\ell-1)R,2dR),\\
 (\Gamma(I_i,I_j))_{I=(R,W,F)}=
 \begin{pmatrix}
 4R&8W&8F\\8W&4RW&4RF\\8F&4RF&4R(dF+abW)
 \end{pmatrix}.
 \end{gathered}
\end{equation}
Equation \eqref{eq:orthq-minimality} follows from
\eqref{bh:minimality}. The smallest-algebra and independence assertions
follow from the argument after \eqref{bh:reduced-table}.
In this case $R$ is already obtained from $F$ by one Laplacian:
$\ell$ is even because $J_1^2=-\Id$, whereas $2m+1$ is odd, so
$d\ne0$ and $R=\Delta F/(2d)$. In particular $F$ is nonharmonic.

For the remaining assertions, fix $x\ne0$ and decompose
\[
 \R^\ell=\R x\oplus U_x^+\oplus U_x^-,\qquad
 U_x^+=\operatorname{span}\{J_1x,\ldots,J_mx\}.
\]
For $y=y_0+y_++y_-$, the exact fibre formulas are
\begin{equation}\label{eq:orthq-fibre}
 W=|x|^2(|y_+|^2+|y_-|^2),\qquad
 F=|x|^2(a|y_+|^2-b|y_-|^2).
\end{equation}
They exhibit regular zeros and both phases. Writing
$F=y^{\mathsf T}M(x)y$, the matrix $M(x)$ has rank $\ell-1\ge4$
for every real $x\ne0$, with kernel $\R x$. It has the same rank
over $\R(x)$, so this quadratic in $y$ cannot factor into linear
forms over that field. Its coefficients, which are homogeneous
quadratics in $x$, have no nonconstant common divisor. A linear
common divisor would give $M(x)=0$ at a nonzero point of a real
hyperplane. A quadratic common divisor would give $M(x)=p(x)B$
for a constant matrix $B$, and $M(x)x=0$ would force $B=0$.
Both contradict the rank assertion. Gauss's lemma therefore proves
real irreducibility.

A constant annihilating direction $(v_1,v_2)$ would give two
vanishing derivatives of distinct bidegrees $(1,2)$ and $(2,1)$.
The second implies $M(x)v_2=0$ for every $x\ne0$, hence $v_2=0$;
interchanging $x,y$ gives $v_1=0$. Thus the ambient space is essential.
Finally $dF+abW=a^2Q+b^2(W-Q)$, so \eqref{eq:orthq-minimality}
identifies the critical set with $W=0$. Equality in real
Cauchy--Schwarz gives \eqref{eq:orthq-singular}. The charts $y=tx$
with $x\ne0$, and their counterparts with $y\ne0$, give its smoothness
and dimension. Smooth local frames in \eqref{eq:orthq-fibre} give
the transverse quadratic cone. At $(x,tx)$ an isometric normal
coordinate is $w=\Pi_{x^\perp}(\delta y-t\delta x)/\sqrt{1+t^2}$; the
nonzero Hessian eigenvalues are $2aR$ and $-2bR$, with multiplicities
$m$ and $\ell-m-1$. Their sum is at least four, so the zero set
cannot be smooth there. This also proves that these critical points
are geometric singularities, as in \cref{lem:quartic-table-geometry}.
\end{proof}

\begin{proof}[Proof of \cref{thm:orthq-infinite-families}]
The weights are $(2,m,4k-m-1)$, and their active sum is $h=4k-1$.
For $k\ge3$, this is at least eleven. The normalization
$T=W$, $F_{\rm gen}=F_{k,m}/(4k-3)$ therefore gives
\eqref{eq:orthq-potential} from \eqref{bh:potentials}, with $W>0$
off the rank-one singular set. The strict-minimality and full-support
conclusions follow from \cref{thm:quartic-table-existence}. In this normalization,
\begin{equation}\label{eq:orthq-potential-defect}
 \begin{split}
 \cL(WF)=8WF\bigl\{&[3(\ell-1)R^2-12W]F^2\\
 &+d[(\ell-5)R^2-4W]WF
 +ab[(\ell-9)R^2-4W]W^2\bigr\},
 \end{split}
\end{equation}
which is \eqref{bh:G8-defect} after the same constant rescaling.
The Hardy coefficient in \eqref{eq:bihomogeneous-Hardy}, with
$N=8k$ and $h=4k-1$, is exactly
\[
 \frac{(8k-17)^2}{4}+8(4k-1)-72
 =16k^2-36k-\frac{31}{4}>0.
\]

At $k=2$, the transverse tangent cones from
\cref{prop:orthq-construction} have signatures $(2,5)$ and $(3,4)$
in $\R^7$. Each is a six-dimensional Lawson cone whose link has
$|A|^2=5$, so its radial Hardy margin is
$(6-2)^2/4-5=-1$. \Cref{lem:tangent-instability-transfer},
applied to the flat $\R^9$
factor and the smooth local model \eqref{eq:orthq-fibre}, transfers
a compactly supported negative variation to the regular cone.
Thus both $k=2$ cones are unstable and not area minimizing.
\end{proof}

\begin{proposition}[Distinction from the preceding families]
\label{prop:orthq-distinct}
For $k\ge3$, the cones in \eqref{eq:orthq-two-families} are not
orthogonally congruent to the standard OT--FKM splitting cones or
to pure full Clifford pullbacks of nondegenerate Lawson quadratic
cones. For each fixed $k$, the two cones are not congruent to one another.
\end{proposition}

\begin{proof}
The splitting cones in $\R^{2\ell}$ have singular set the two
coordinate $\ell$-planes, whereas \eqref{eq:orthq-singular} has
dimension $\ell+1$. For a full Clifford quadratic map
$P:\R^{2\ell}\to\R^p$, the identity
$DP\,DP^{\mathsf T}=4R\Id_p$ shows that the nonvertex singular
set of a nondegenerate Lawson pullback is $P^{-1}(0)$, of codimension
$p$ if nonempty. For the balanced target this is also recorded in
\cite[Remark~5.2.6(1)]{RiedlerThesis25}. Congruence would force $p=\ell-1$.
The Clifford module restriction
\cite[Section~3.5]{FKM81}, equivalently
\cite[Table~4.1]{RiedlerThesis25}, would then require \(\delta(\ell-2)=2^{2k-1}\mid4k=\ell\),
which is impossible for $k\ge3$ because $2^{2k-1}>4k$.
Finally the transverse signatures $(2,4k-3)$ and $(3,4k-4)$
are distinct even after interchanging their two entries, and hence
distinguish the two sequences.
\end{proof}

\begin{remark}[The weights $(2,5,10)$]
\label{ex:quartic-2510}
Taking $\ell=16$ and $m=5$ in \cref{prop:orthq-construction}
gives the weights $(2,5,10)$ and $F=13Q-4W$ on $\R^{32}$.
The required Clifford system exists since $\delta(6)=8\mid16$.
Here $h=15$, so \cref{thm:quartic-table-existence} gives the
potential $WF$, strict area minimization, and strict stability.
\end{remark}
